\section{Literature Review}\label{literature}
Collusion in public procurement auctions has been extensively studied through the lenses of game theory, empirical analysis, and policy evaluation. A rich theoretical literature explains how bidders can sustain collusive agreements in repeated auction environments. Key mechanisms include bid rotation, where cartel members take turns winning contracts, and strategies to deter outsider entry. For example, dynamic models show that bidding rings can allocate wins intertemporally without side payments by rotating bids in a predetermined order, ensuring that each cartel member secures a share of contracts over time. Auction format and information flow also matter: tacit coordination can emerge even without explicit communication, especially under auction designs that facilitate monitoring of rivals’ behavior. Notably, uniform-price auctions tend to relax competitive pressure more than pay-as-bid formats, thereby facilitating implicit collusion \citep{porter1993detecting}.

Empirical evidence since the 2000s demonstrates that collusive schemes are not merely theoretical constructs but real phenomena observed across different regions. Numerous documented cartels in procurement markets illustrate various collusive strategies. In Japan, for instance, construction contractors have systematically rigged bids on public works projects, with nationwide tender data revealing widespread coordination \citep{kawai2023using}. Within such schemes, firms often allocate subsequent contracts to the member that has gone the longest without winning, a clear indicator of bid rotation. Similar patterns have been identified in Italy, Scandinavian countries, and in Brazil’s public contracting, where studies document that collusive rings inflate contract prices and coordinate bidding behavior through predictable bid rotations and artificially close bids \citep{bajari2003detecting}. These patterns persist over time, leading to higher procurement costs and diminished competition \citep{conley2016detecting}.

Given the covert nature of collusion, researchers have focused on developing robust detection methods. Traditional econometric approaches leverage the insight that collusion distorts the statistical properties of bids. For example, tests have been developed to assess whether bidders’ cost estimates or bid rankings are exchangeable under competitive conditions, with systematic deviations serving as red flags for collusion \citep{bajari2003detecting}. More recent work has advanced these methods by applying Regression Discontinuity (RD) designs that exploit near-tie auctions, where the winning margin is nearly zero. Under true competition, firms in such close contests should display similar characteristics, and any systematic differences can indicate collusive coordination \citep{imhof2018screening}. These techniques have successfully detected collusion in various procurement settings, highlighting that even sophisticated cartels leave measurable traces in bidding behavior.

Policy responses to collusion have evolved significantly since 2000. Regulatory frameworks now emphasize proactive detection and robust enforcement. Leniency programs, which incentivize cartel members to self-report by offering reduced penalties, have been widely implemented and have been shown to destabilize collusive arrangements \citep{oecd2015}. In parallel, governments have refined auction design—by increasing transparency, enforcing strict bid submission protocols, and adopting electronic procurement platforms—to reduce the opportunities for bid rigging. Empirical studies indicate that such measures, when effectively enforced, lead to lower procurement costs and more competitive outcomes \citep{bosio2022public}.

Looking ahead, recent innovations in machine learning and network analysis are poised to further enhance collusion detection. These advanced methods can process vast quantities of bidding data to identify subtle patterns that traditional econometric screens may overlook, offering a more comprehensive toolkit for uncovering collusive behavior. As digital procurement continues to expand globally, integrating these tools into routine oversight could significantly improve market competitiveness and safeguard public resources.

In sum, the post-2000 literature paints a comprehensive picture of collusion in procurement: it is grounded in robust theoretical frameworks, confirmed by empirical evidence across diverse jurisdictions, and addressed by increasingly sophisticated detection and regulatory strategies. This body of work not only deepens our understanding of collusive behavior but also provides practical guidance for interventions aimed at preserving competitive procurement markets.
